\documentclass[a4paper,11pt]{article}
\usepackage[T1]{fontenc}
\usepackage[utf8]{inputenc}
\usepackage{amsmath}
\usepackage{siunitx}
\usepackage[margin=2.5cm]{geometry}

\title{ATCG I Sheet R03 Solution}
\author{
\large Kamran Moradnejad (3099229) \\
\large Pranav Yadav (50327698) \\
\large Saransh Jain (50262919)
}
\date{May 7, 2026}

\begin{document}
\maketitle

\section*{Assignment 2: Reflection Equation}
The sun is modeled as a distant disk emitter with constant radiance
\[
L_s = 20.045\,\si{\mega\watt\per\square\meter\per\steradian}
= 2.0045 \times 10^7\,\si{\watt\per\square\meter\per\steradian}.
\]
Its direction encloses an angle of $45^\circ$ with the table normal. Using the solar radius
\[
R_s = 6.9634 \times 10^8\,\si{\meter}
\]
and the mean Earth-sun distance
\[
D = 1.495978707 \times 10^{11}\,\si{\meter},
\]
the angular radius of the sun is
\[
\alpha = \arcsin\left(\frac{R_s}{D}\right),
\]
and the solid angle of the solar disk is
\[
\Omega_s = 2\pi(1-\cos\alpha) \approx 6.8068 \times 10^{-5}\,\si{\steradian}.
\]

Therefore the irradiance on the table is
\[
E_{\text{table}} = \int_{\Omega_s} L_s \cos 45^\circ \, d\omega
= L_s \cos 45^\circ \, \Omega_s
\approx 964.8\,\si{\watt\per\square\meter}.
\]

The table is perfectly diffuse with albedo $\rho = 0.5$, hence its outgoing radiance is
\[
L_{\text{table}} = \frac{\rho}{\pi} E_{\text{table}}
= \frac{0.5}{\pi} \cdot 964.8
\approx 153.55\,\si{\watt\per\square\meter\per\steradian}.
\]

The lens is a disk with diameter $d = 5\,\si{\centi\meter}$ and radius
\[
r_{\text{lens}} = 0.025\,\si{\meter},
\qquad
A_{\text{lens}} = \pi r_{\text{lens}}^2 \approx 1.9635 \times 10^{-3}\,\si{\square\meter}.
\]
The table is a square of side length $0.8\,\si{\meter}$, so with half side length $a = 0.4\,\si{\meter}$ and lens distance $h = 1.2\,\si{\meter}$, the irradiance at the lens is
\[
E_{\text{lens}}
= \int_A L_{\text{table}} \frac{\cos\theta_t \cos\theta_\ell}{r^2}\, dA
= L_{\text{table}} \int_{-a}^{a}\int_{-a}^{a}
\frac{h^2}{(x^2+y^2+h^2)^2}\, dy\, dx.
\]
Numerically,
\[
\int_{-0.4}^{0.4}\int_{-0.4}^{0.4}
\frac{1.2^2}{(x^2+y^2+1.2^2)^2}\, dy\, dx
\approx 0.38741,
\]
so
\[
E_{\text{lens}} \approx 153.55 \cdot 0.38741
\approx 59.49\,\si{\watt\per\square\meter}.
\]

The radiant power received by the lens is therefore
\[
\Phi_{\text{lens}} = E_{\text{lens}} A_{\text{lens}}
\approx 59.49 \cdot 1.9635 \times 10^{-3}
\approx 1.168 \times 10^{-1}\,\si{\watt}.
\]

Hence,
\[
\boxed{E_{\text{lens}} \approx 59.49\,\si{\watt\per\square\meter}}
\]
and
\[
\boxed{\Phi_{\text{lens}} \approx 0.1168\,\si{\watt}.}
\]

\section*{Assignment 3: Neumann Series Rendering Equation}
\textbf{a)} In
\[
L = \sum_{i=0}^{\infty} T^i L_e,
\]
the quantities are:
\begin{itemize}
    \item $L$: the unknown radiance distribution in the scene,
    \item $L_e$: the emitted radiance distribution,
    \item $T$: the transport operator that maps incident radiance to reflected radiance after one scattering event.
\end{itemize}

\textbf{b)} The first summands have the following meaning:
\begin{itemize}
    \item $T^0 L_e = L_e$: directly emitted radiance without any bounce,
    \item $T^1 L_e$: radiance after exactly one bounce,
    \item $T^2 L_e$: radiance after exactly two bounces.
\end{itemize}

\textbf{c)} If the Neumann series converges, then
\[
\sum_{i=0}^{\infty} T^i = (I-T)^{-1},
\]
so the rendering equation can be written as
\[
\boxed{L = (I-T)^{-1} L_e.}
\]

\end{document}
